\section{The \texorpdfstring{\toolname}{PhishProof} Detector}
\label{sec:method}

\toolname{} turns the Grounded Evidence-Agreement score of \eqref{eq:gea} into a
calibrated, selective verdict through the four components $C_1$--$C_4$ previewed in
\autoref{sec:form-overview} and drawn in \autoref{fig:overview}. The four subsections
below develop them in the order a page flows through the system -- elicit evidence,
agree on it, ground it, then decide -- after which we give the end-to-end procedure,
its guarantee, and its cost.

\subsection{Evidence agents and the shared cue set}
\label{sec:elicit}
This component produces the per-agent outputs of \eqref{eq:output}. Each member of the
panel is an \emph{evidence agent}: prompted with the page, it \emph{plans} which
signals and modalities to examine, calls the verification tools to obtain them, and
returns a verdict together with the grounded cues it found, drawn from the fixed schema
$\Uschema$ and emitted by constrained decoding as well-typed $(t,v)$ pairs rather than
free text (\req{1}). Constraining the output to the schema is what makes the cue
auditable: it is a typed claim a tool can later re-derive, not prose to be
re-interpreted. The schema is small and auditable -- the brand
claim, the form-action domain, credential intent, the logo brand, a brand--domain
consistency cue derived from the agreed brand and the host, and, where the capture
provides them, the certificate organization and redirect target (both absent in the
corpus of \autoref{sec:experiments}, hence inert there) -- and is shared by every agent, so their cues
are directly comparable. The cues are pooled into the page's \emph{shared cue set},
where each cue is typed, tool-checkable, and tagged with the agent that asserted it, so
the verdict rests on a list of cited, verifiable reasons -- a brand claim paired with
off-brand infrastructure -- rather than an opaque score. Agent \emph{diversity} is a
design choice, not an afterthought: the three agents come from different model families
and at least one is a vision-language model that reads the screenshot, so that
\emph{agreement} on a cue reflects independent corroboration rather than a shared prior
the agents happen to hold (its value is isolated in \autoref{sec:experiments}).

\subsection{Evidence agreement}
\label{sec:agree}
This component scores the shared cue set of \autoref{sec:elicit}. To \emph{aggregate}
and compare cues across agents, two agents must be able to assert the ``same'' cue
despite surface differences, so we first \emph{normalize values}: domains are reduced
to their registrable eTLD+1 and brands to a canonical lexicon entry, so that
``\textsf{paypal-secure.com/login}'' and ``\textsf{paypal-secure.com/signin}'' count as
one form-action cue. On the normalized cues, the consensus set $\cons(\page)$ collects
those a strict majority of the agents assert, and the evidence-level agreement
$\agr(\page)$ of \eqref{eq:agree} is the share of all cues raised that fall in this
consensus (\req{2}). One refinement keeps the rule fair to single-source evidence: a
\emph{perceptual} cue visible to only one modality -- the logo, which only the
vision agent can read -- is admitted to the consensus on its own, since it cannot reach a
majority when a lone agent sees the screenshot; every other type still requires a strict
majority. Agreement is therefore high only when the agents converge on the
\emph{same} evidence, and -- unlike a label-vote concordance -- is unmoved by their
merely concurring on the verdict.

\subsection{Tool-grounding}
\label{sec:ground}
This component verifies the consensus cues from \autoref{sec:agree} against the page.
Agreement alone cannot catch a \emph{shared} hallucination -- a cue the whole panel
asserts but the page does not carry -- so each consensus cue is re-derived from the
page by the tool for its type, $\gver(u,\page)$ (\req{3}). The structural tools verify
deterministically and return $0$ or $1$: the DOM is parsed for the form-action domain,
the certificate for its organization, and the redirect chain for its final target. The
perceptual tool (\textsf{logo-brand}) matches the rendered logo to the claimed brand
and returns a similarity in $[0,1]$. The brand claim is grounded by a check that the page
itself presents the claimed brand; this tool slot can equally be filled by an existing
specialized detector \emph{used as a tool} -- e.g.\ a reference-based visual
matcher~\cite{lin2021phishpedia} or, in principle, an LLM brand-and-domain
checker~\cite{tan2024phishllm}. We verify substitution empirically for the Phishpedia (D1)
matcher: replacing the page-content brand check leaves the selective ranking essentially
unchanged (\autoref{sec:experiments}), so the panel can \emph{subsume} that detector's
signal as a grounding tool rather than re-derive it -- a property we measure for D1 rather
than assume for every specialized detector. The groundedness $\grd(\page)$ of
\eqref{eq:ground} averages these tool outputs over the consensus, and through the
product $\geascore=\agr\cdot\grd$ of \eqref{eq:gea} the score collapses whenever the
agreed evidence does not survive verification. On a clean corpus the agreed cues nearly
always survive that check, so $\grd$ is close to constant and contributes little to the
\emph{ranking} (\autoref{sec:exp_rq3}); its job is to keep the score honest in the
regime the clean corpus lacks, where an attacker corrupts the very cue the panel cites.
There $\grd$ falls, $\geascore$ collapses, and the detector abstains rather than acts on
corrupted evidence -- the fail-safe behaviour we measure directly in
\autoref{sec:exp_adv}.

\subsection{Calibrated selection}
\label{sec:select}
Finally, \toolname{} turns the score $\geascore(\page)$ into an act/abstain decision.
Because phishing labels are available at collection time -- a flagged URL is confirmed
by the public feed it is drawn from~\cite{openphish,phishtank} -- we can set the
operating point on a small held-out labeled split rather than guess it. On that split
we fit a monotonic (isotonic) calibrator $\kappa:\geascore\mapsto\hat\Prob(\hat y = y)$,
which is order-preserving and so leaves the score's ranking -- hence the
risk--coverage curve -- intact while turning it into a probability of correctness, and
we pick the smallest threshold $\thr$ that meets a target selective risk. The detector
then acts with $\hat y(\page)$ when $\geascore(\page)\ge\thr$ and abstains otherwise,
escalating the page to deeper analysis or a human. When it acts, it returns the
consensus cue set $\cons(\page)$ as a verifiable explanation.

\subsection{Procedure, guarantee, and cost}
\label{sec:algo}
\begin{algorithm}[!h]
\caption{\toolname{} inference on a page $\page$}
\label{alg:gea}
\begin{algorithmic}[1]
\State \textbf{for} $m_i\in\panel$: $\ (y_i,\evset_i)\gets m_i(\page)$
       \hfill\Comment{$C_1$ agents call tools, \eqref{eq:output}}
\State $\hat y\gets \mathrm{maj}_i\,y_i$;\ \
       $E\gets\bigcup_i\evset_i$ \hfill\Comment{shared cue set}
\State $\cons\gets\{u: u$ asserted by $>\!M/2$ agents, value-normalized$\}$
\State $\agr\gets \lvert\cons\rvert\,/\,\lvert E\rvert$
       \hfill\Comment{$C_2$, \eqref{eq:agree}}
\State \textbf{for} $u\in\cons$: compute $\gver(u,\page)$;\ \
       $\grd\gets \mathrm{mean}_{u\in\cons}\gver(u,\page)$
       \hfill\Comment{$C_3$, \eqref{eq:ground}}
\State $\geascore\gets \agr\cdot\grd$;\ \ $\hat\rho\gets\kappa(\geascore)$
       \hfill\Comment{\eqref{eq:gea}}
\State \textbf{if} $\geascore\ge\thr$ \textbf{return} $(\hat y,\hat\rho,\cons)$;
       \textbf{else return} $\bot$
\end{algorithmic}
\end{algorithm}

\sstitle{Walk-through}
\autoref{alg:gea} runs the agent panel once, each agent calling its tools to cite its
cues (line~1), takes the majority verdict (line~2), and then does the work that
distinguishes \toolname{} from label-level reliability: it forms the consensus set
(line~3), scores agreement on that evidence rather than on the verdict (line~4),
checks each agreed cue against the page with its tool so a shared hallucination cannot
inflate the score (line~5), and combines the two factors into $\geascore$ (line~6).
Calibration and the selective decision (line~7) then map the score to a trust value and
apply the threshold fitted on the labeled split. The novelty is concentrated in
lines~3--6: the quantity that gates the verdict is grounded agreement on the
\emph{cues}, not concordance of labels, and the returned consensus set is
simultaneously the trust signal and the explanation -- one object serving both
\req{2}--\req{3} and the audit trail of \req{1}.

\sstitle{Guarantee}
Tool-grounding is the component that rules out a shared hallucination inflating the
score, and under one assumption on the tools we can state this precisely.
\begin{assumption}[Tool soundness]\label{ass:sound}
For each cue type, the tool $\gver(u,\page)=1$ if $u$ holds on $\page$ and $0$
otherwise (perceptual types: $\gver(u,\page)$ is the true match score).
\end{assumption}
\begin{proposition}\label{prop:hall}
Under \autoref{ass:sound}, any consensus cue that does not hold on $\page$
contributes $0$ to $\grd(\page)$; hence $\geascore(\page)$ is at most the grounded
fraction of the consensus, and $\geascore(\page)=0$ whenever the entire consensus
is hallucinated.
\end{proposition}
\begin{proof}[Proof sketch]
$\grd$ is the mean of $\gver$ over $\cons$ (\eqref{eq:ground}); by
\autoref{ass:sound} a non-holding cue adds $0$ to that mean, so
$\grd\le\lvert\{u\in\cons: u\text{ holds}\}\rvert/\lvert\cons\rvert$, and
$\geascore=\agr\cdot\grd\le\grd$ inherits the bound. If no consensus cue holds,
$\grd=0$ and thus $\geascore=0$.
\end{proof}
The proposition is modest by design: it bounds, but does not eliminate, gauge error,
since a tool is sound for \emph{what} it checks yet a page can still carry
true-yet-misleading evidence the tool confirms; accordingly we report measured tool
precision and recall against gold in \autoref{sec:experiments} rather than rely on
\autoref{ass:sound} as if it held exactly.

\sstitle{Complexity}
Per page, the cost has two parts: the $M$ agent calls of line~1 -- each a model call
plus the tool calls it issues -- and the verification of the $\lvert\cons\rvert$
consensus cues in line~5. The aggregation steps (lines~2--4,~6) are linear in the
number of cues raised and negligible beside a model call. Each structural tool is an
$\bO(1)$ DOM, certificate, or redirect lookup and the single perceptual check is one
embedding comparison, so verification is $\bO(\lvert\cons\rvert)$ and the per-page cost
is dominated by the $M{=}3$ agent calls; it is independent of the schema size beyond
the cues actually asserted, and there is no training. Relative to the strongest
non-agentic baseline -- a single-model detector -- this is a fixed $M{\times}$ inference
overhead. Selective operation amortizes it: only the abstaining fraction
$1-\phi(\thr)$ is escalated to deeper analysis, so the expected per-page budget is tuned
by the very threshold $\thr$ that sets the risk--coverage operating point
(\autoref{sec:experiments}).
